\documentclass[12pt]{article}
\usepackage[utf8]{inputenc}
\usepackage[T1]{fontenc}
\usepackage{hyperref}
\title{Tilquiapan Zapotec Language Resources for K-12 Curriculum Development}
\author{Compiled by Claude Sonnet 5.5 (claude-sonnet-5-5) for Andr\'as Kornai}
\date{}
\begin{document}
\maketitle

\section{Introduction}
Tilquiapan Zapotec (ISO 639-3 zts, Glottocode tilq1235) is a Zapotecan language of southern Oaxaca, Mexico, with about 5,000 speakers. It is also called Zapoteco de San Miguel Tilquiapam. It is close to Santa In\'es Yatzechi Zapotec. A practical orthography was developed with speakers between 2000 and 2008. The searches found few resources, and none designed for schools.

\section{Primary Educational Resources}
No primers, textbooks or school curricula were found in the searches.

\section{Linguistic and Reference Materials}
\begin{itemize}
\item Wikipedia overview (phonology, grammar, writing): \url{https://en.wikipedia.org/wiki/Tilquiapan_Zapotec}
\item PHOIBLE phonological inventory: \url{https://phoible.org/languages/zts}
\item Academic profile of Beth Merrill, whose work describes the language: \url{https://independent.academia.edu/BethMerrill2}
\end{itemize}

\section{Religious and Community Resources}
\begin{itemize}
\item Scripture Earth listing: \url{https://www.scriptureearth.org/00eng.php?iso=zts}
\item Joshua Project profile: \url{https://joshuaproject.net/people_groups/14711}
\end{itemize}

\section{Digital and Audio Resources}
\begin{itemize}
\item Global Recordings Network audio: \url{https://globalrecordings.net/en/language/zts}
\end{itemize}

\section{Curriculum Development Context}
The language is written in Latin script with an orthography worked out with speakers. Documentation is mostly linguistic, and the Bible and recorded lessons are the only applied materials found. A curriculum would need new materials developed with the community.
\end{document}